% articulo-11-citas.tex - Citas exentas y separador entre bloques.
% (quoting, \quotingbreak)
%
% Chicago § 2.23, RAE § T-74.
% indentfirst=false: el primer párrafo de la cita no lleva sangría (norma).
% A partir del segundo párrafo sí lleva sangría: se restaura con \everypar.
% Sangrado solo por la izquierda (Chicago, por defecto).
%
% Variables YAML:
%   (sin declarar)                  →  sangrado solo por la izquierda (por defecto)
%   citas-sangrado-simetrico: true  →  sangrado simétrico izquierda y derecha
%
% \quotingbreak: separador entre bloques de cita independientes fusionados en
% un único entorno quoting por el filtro fusionar-citas.lua.
% Dos citas Markdown separadas por línea en blanco fuera del «>» generan dos
% BlockQuote independientes; el filtro los fusiona con \quotingbreak entre
% ellos. \par cierra el párrafo en curso; \noindent suprime la sangría del
% siguiente párrafo (norma de inicio de cita).
%
% El filtro distingue dos casos:
%   > Párrafo uno        →  línea en blanco DENTRO del ">": un solo entorno
%   >                        quoting con dos párrafos; el segundo lleva sangría.
%   > Párrafo dos
%
%   > Párrafo uno        →  línea en blanco FUERA del ">": dos entornos
%                            quoting independientes fusionados con \quotingbreak;
%   > Párrafo dos            ninguno lleva sangría en primera línea.

\usepackage{quoting}
\quotingsetup{
  font        = small,
  indentfirst = false,
  leftmargin  = \parindent,
  rightmargin = $if(citas-sangrado-simetrico)$\parindent$else$0pt$endif$,
  vskip       = 0pt             % sin espacio extra antes/después de la cita
}

% Restaurar sangría a partir del segundo párrafo dentro de quoting.
% Solo tiene efecto cuando hay varios párrafos en el mismo bloque de cita.
\AtBeginEnvironment{quoting}{%
  \everypar{\parindent=\dimexpr\parindent\relax}%
}

\newcommand{\quotingbreak}{\par\noindent}
